\subsubsection{Carath\'eodory--Toeplitz 语义下的端点热核探针：Toeplitz 主子矩阵实现、单调极限与可审计推论链}\label{subsubsec:xi-endpoint-heat-kernel-probe-toeplitz}
本段将端点热核探针置于 Carath\'eodory--Toeplitz 语义之内：探针既可表述为 \(\TT\) 上正测度的端点权重矩（从而适配端点递归与局部喷流分析），亦可在有限阶 Toeplitz 主子矩阵上由单一二次型给出（从而适配矩证书/可见证书的离线复核与复杂度评估）。

\paragraph{Toeplitz 实现与严格恒等式}
沿用定理 \ref{thm:xi-endpoint-heat-kernel-probe-third-order-expansion} 中端点热核探针的记号 \(a_N=\int_{\TT}\bigl|\frac{1-\xi}{2}\bigr|^{2N}\,\dd\mu(\xi)\)。
令 \(\mu\) 为 \(\TT\) 上有限正 Borel 测度，并令其 Fourier 系数
\[
c_h:=\int_{\TT}\xi^{-h}\,\dd\mu(\xi),\qquad h\in\ZZ,
\]
从而 \(c_{-h}=\overline{c_h}\) 且 Toeplitz 主子矩阵
\[
\mathbf{T}_N:=(c_{j-k})_{0\le j,k\le N}\in\Mat_{N+1}(\CC)_{\mathrm{Herm}}
\]
半正定（Carath\'eodory--Toeplitz 完备性，参见附录相应刻画与 \S\ref{subsubsec:xi-pontryagin-toeplitz-blaschke-tomography-tate} 的记号约定）。
定义端点热核探针向量，并引入其 Toeplitz 二次型表示
\begin{equation}\label{eq:xi-endpoint-heat-probe-bN-aN}
\boxed{\;
b^{(N)}_k:=(-1)^k\binom{N}{k}\ (0\le k\le N),\qquad
a_N=4^{-N}\,(b^{(N)})^\ast\,\mathbf{T}_N\,b^{(N)}\;}
\end{equation}
并记端点压缩函数
\begin{equation}\label{eq:xi-endpoint-rxi-def}
\boxed{\;
r(\xi):=\frac{|1-\xi|^2}{4}\in[0,1]\qquad(\xi\in\TT).\;}
\end{equation}

\begin{proposition}[端点热核探针的 Toeplitz--Herglotz 同一]\label{prop:xi-endpoint-heat-probe-toeplitz-herglotz-identity}
在上述记号下，对所有 \(N\ge 0\) 成立严格恒等式
\begin{equation}\label{eq:xi-endpoint-heat-probe-identity}
\boxed{\;
a_N
=\frac{1}{4^N}\int_{\TT}|1-\xi|^{2N}\,\dd\mu(\xi)
=\int_{\TT}r(\xi)^N\,\dd\mu(\xi)\ \ge\ 0.\;}
\end{equation}
\end{proposition}
\begin{proof}
由定义与 Fubini 交换求和积分，
\[
(b^{(N)})^\ast\mathbf{T}_N b^{(N)}
=\sum_{j,k=0}^{N}\overline{b_j^{(N)}}\,b_k^{(N)}\,c_{j-k}
=\int_{\TT}\sum_{j,k=0}^{N}\overline{b_j^{(N)}}\,b_k^{(N)}\,\xi^{-(j-k)}\,\dd\mu(\xi).
\]
注意到
\(
\sum_{k=0}^{N}b_k^{(N)}\xi^{k}=(1-\xi)^N
\)
以及
\(
\sum_{j=0}^{N}\overline{b_j^{(N)}}\xi^{-j}=(1-\xi^{-1})^N
=\overline{(1-\xi)^N}
\)
（在 \(\TT\) 上 \( \xi^{-1}=\overline{\xi}\)）。
故 integrand 化简为 \(|1-\xi|^{2N}\)，得到第一等号；第二等号由 \eqref{eq:xi-endpoint-rxi-def} 直接改写。
最后由 \(\mu\) 正性得 \(a_N\ge 0\)。证毕。
\end{proof}

\paragraph{单调极限、谱隙指数界与端点质量}
记端点原子质量 \(m_{-1}:=\mu(\{-1\})\)。

\begin{corollary}[端点质量的单调读出与谱隙指数界]\label{cor:xi-endpoint-heat-probe-monotone-limit-gap}
对 \(a_N\) 有单调性与极限
\begin{equation}\label{eq:xi-endpoint-heat-probe-monotone-limit}
\boxed{\;
0\le a_{N+1}\le a_N,\qquad
a_N\downarrow m_{-1}\quad(N\to\infty).\;}
\end{equation}
并且若存在常数 \(r_\star\in[0,1)\) 使得
\[
\sup_{\xi\in\operatorname{supp}(\mu)\setminus\{-1\}}r(\xi)\le r_\star,
\]
则对所有 \(N\ge 0\) 有指数型误差界
\begin{equation}\label{eq:xi-endpoint-heat-probe-gap-exp-bound}
\boxed{\;
0\le a_N-m_{-1}\le r_\star^{\,N}\,\mu(\TT\setminus\{-1\}).\;}
\end{equation}
\end{corollary}
\begin{proof}
由 \eqref{eq:xi-endpoint-heat-probe-identity}，\(a_N=\int r^N\,\dd\mu\) 且 \(0\le r\le 1\)。
故 \(r^{N+1}\le r^N\) 给出单调性。
又因 \(r(\xi)=1\) 当且仅当 \(\xi=-1\)，由单调收敛定理得
\(
\lim_{N\to\infty}a_N=\mu(\{r=1\})=\mu(\{-1\})=m_{-1}
\)。
指数界由
\(
a_N-m_{-1}=\int_{\TT\setminus\{-1\}}r(\xi)^N\,\dd\mu(\xi)\le r_\star^N\mu(\TT\setminus\{-1\})
\)
得到。证毕。
\end{proof}

\paragraph{谱隙半径的极限刻画与内部校验约束}
定义端点外支撑的谱隙半径
\begin{equation}\label{eq:xi-endpoint-heat-probe-rstar-def}
\boxed{\;
r_\ast:=\sup_{\xi\in\operatorname{supp}(\mu)\setminus\{-1\}}r(\xi)\in[0,1].\;}
\end{equation}

\begin{theorem}[指数根与比值极限决定谱隙半径]\label{thm:xi-endpoint-heat-probe-rstar-root-ratio}
若 \(m_{-1}<\mu(\TT)\)（即 \(\mu(\TT\setminus\{-1\})>0\)），则
\begin{equation}\label{eq:xi-endpoint-heat-probe-root-limit}
\boxed{\;\lim_{N\to\infty}(a_N-m_{-1})^{1/N}=r_\ast.\;}
\end{equation}
进一步地，令
\(
\rho_N:=\frac{a_{N+1}-m_{-1}}{a_N-m_{-1}}
\)
（\(N\) 充分大时分母为正），则 \(\rho_N\) 单调递增且
\begin{equation}\label{eq:xi-endpoint-heat-probe-ratio-limit}
\boxed{\;\rho_N\uparrow r_\ast.\;}
\end{equation}
\end{theorem}
\begin{proof}
由 \eqref{eq:xi-endpoint-heat-probe-identity}，
\(
a_N-m_{-1}=\int_{\TT\setminus\{-1\}} r(\xi)^N\,\dd\mu(\xi)
\)。
上界由 \(r(\xi)\le r_\ast\) 给出
\(
(a_N-m_{-1})^{1/N}\le r_\ast\,\mu(\TT\setminus\{-1\})^{1/N}\to r_\ast
\)。
对下界，任取 \(\varepsilon>0\)。由 \(r_\ast\) 为支撑上的上确界，存在 \(\xi_0\in\operatorname{supp}(\mu)\setminus\{-1\}\) 使 \(r(\xi_0)>r_\ast-\varepsilon\)。
由支撑定义取邻域 \(U\ni\xi_0\) 使 \(\mu(U)>0\) 且 \(r(\xi)\ge r_\ast-\varepsilon\) 于 \(U\) 成立。
于是
\(
a_N-m_{-1}\ge \int_U r(\xi)^N\,\dd\mu(\xi)\ge (r_\ast-\varepsilon)^N\mu(U)
\)，从而 \(\liminf (a_N-m_{-1})^{1/N}\ge r_\ast-\varepsilon\)。
令 \(\varepsilon\downarrow 0\) 得 \eqref{eq:xi-endpoint-heat-probe-root-limit}。
\par
对比值 \(\rho_N\)，先由 \eqref{eq:xi-endpoint-heat-probe-logconvex}（见后述命题）得 \(\rho_N\) 单调递增。
另一方面，由
\[
\rho_N=\frac{\int_{\TT\setminus\{-1\}} r(\xi)^{N+1}\,\dd\mu(\xi)}{\int_{\TT\setminus\{-1\}} r(\xi)^{N}\,\dd\mu(\xi)}
\]
可知 \(\rho_N\le r_\ast\)（加权平均界）。
为证极限下界，任取 \(\varepsilon>0\)。令 \(A_{\varepsilon}:=\{\xi\in\TT\setminus\{-1\}: r(\xi)\ge r_\ast-\varepsilon\}\)。
由 \(r_\ast\) 为支撑上确界可取 \(\varepsilon/2\) 使 \(\mu(A_{\varepsilon/2})>0\)。
则
\[
1-\frac{\int_{A_{\varepsilon}} r^N\,\dd\mu}{\int r^N\,\dd\mu}
=\frac{\int_{\TT\setminus(\{-1\}\cup A_{\varepsilon})} r^N\,\dd\mu}{\int r^N\,\dd\mu}
\le
\frac{(r_\ast-\varepsilon)^N\,\mu(\TT)}{(r_\ast-\varepsilon/2)^N\,\mu(A_{\varepsilon/2})}\xrightarrow[N\to\infty]{}0.
\]
于是
\[
\rho_N\ge (r_\ast-\varepsilon)\cdot \frac{\int_{A_{\varepsilon}} r^N\,\dd\mu}{\int r^N\,\dd\mu}=(r_\ast-\varepsilon)\,(1+o(1)).
\]
令 \(\varepsilon\downarrow 0\) 并结合 \(\rho_N\le r_\ast\) 与单调性，得 \(\rho_N\uparrow r_\ast\)。证毕。
\end{proof}

\begin{proposition}[完全单调性、对数凸性与差分非负]\label{prop:xi-endpoint-heat-probe-complete-monotone-logconvex}
令 \(s_N:=a_N-m_{-1}=\int_{\TT\setminus\{-1\}}r(\xi)^N\,\dd\mu(\xi)\)。
则对任意整数 \(k\ge 0\)、\(N\ge 0\) 有完全单调性
\begin{equation}\label{eq:xi-endpoint-heat-probe-complete-monotone}
\boxed{\;(-1)^k\,\Delta^k s_N\ge 0,\;}
\end{equation}
其中 \(\Delta\) 为向前差分算子。
并且 \((s_N)\) 对数凸：
\begin{equation}\label{eq:xi-endpoint-heat-probe-logconvex}
\boxed{\;s_N^2\le s_{N-1}s_{N+1}\qquad(N\ge 1).\;}
\end{equation}
\end{proposition}
\begin{proof}
由 \(s_N=\int r^N\,\dd\mu\) 且 \(0\le r<1\)（\(\TT\setminus\{-1\}\) 上），计算
\[
\Delta^k s_N=\int r^N(r-1)^k\,\dd\mu,
\qquad
(-1)^k\Delta^k s_N=\int r^N(1-r)^k\,\dd\mu\ge 0,
\]
得 \eqref{eq:xi-endpoint-heat-probe-complete-monotone}。
对数凸性由 Cauchy--Schwarz：
\[
s_N^2=\Bigl(\int r^N\,\dd\mu\Bigr)^2
\le \Bigl(\int r^{N-1}\,\dd\mu\Bigr)\Bigl(\int r^{N+1}\,\dd\mu\Bigr)=s_{N-1}s_{N+1}.
\]
证毕。
\end{proof}

\paragraph{端点幂律喷流与 dyadic 差分比值的指数读出}
以下讨论 \(r(\xi)\uparrow 1\) 的端点局部几何如何在 \(a_N\) 的多尺度差分中显影。

\begin{theorem}[端点幂律密度诱导的探针主项]\label{thm:xi-endpoint-heat-probe-powerlaw-leading}
设 \(\mu\) 在 \(-1=e^{\mathrm{i}\pi}\) 的邻域内分解为
\[
\mu=m_{-1}\delta_{-1}+f(\theta)\,\dd\theta+\mu_\perp,
\]
其中 \(\mu_\perp\) 在 \(\theta=\pi\) 的某邻域上为零。
假设存在 \(\alpha>-1\) 与常数 \(K_\pm\ge 0\) 使得当 \(\theta\to\pi^\pm\) 时
\[
f(\theta)\sim
\begin{cases}
K_+(\theta-\pi)^\alpha,& \theta\downarrow\pi,\\
K_-(\pi-\theta)^\alpha,& \theta\uparrow\pi.
\end{cases}
\]
记 \(p:=\frac{\alpha+1}{2}>0\)。
则当 \(N\to\infty\) 时有渐近
\begin{equation}\label{eq:xi-endpoint-heat-probe-powerlaw-asymptotic}
\boxed{\;
a_N-m_{-1}
=c\,N^{-p}+o(N^{-p}),\qquad
c:=2^\alpha\,\Gamma(p)\,(K_++K_-)\;}
\end{equation}
且 \(c>0\) 当且仅当 \(K_++K_->0\)。
\end{theorem}
\begin{proof}
由 \eqref{eq:xi-endpoint-heat-probe-identity}，
\(
a_N-m_{-1}=\int_{\TT\setminus\{-1\}}r(\xi)^N\,\dd\mu(\xi)
\)。
在 \(\mu_\perp\) 的支撑上 \(r(\xi)\le q<1\)，其贡献为 \(O(q^N)\)。
对绝对连续部分，令 \(\theta=\pi+\varphi\)，并用
\(
r(e^{\mathrm{i}(\pi+\varphi)})=\cos^2(\varphi/2)
\)。
在 \(|\varphi|\) 小时
\(
\log\cos^2(\varphi/2)=-\varphi^2/4+O(\varphi^4)
\)，因此主贡献来自 \(|\varphi|\asymp N^{-1/2}\)。
对 \(\varphi>0\) 的一侧，代入尺度变换 \(\varphi=2y/\sqrt N\) 并用幂律假设
\(
f(\pi+\varphi)\sim K_+\varphi^\alpha
\) 得
\[
\int_{0}^{\varepsilon}\cos^{2N}\!\left(\frac{\varphi}{2}\right)f(\pi+\varphi)\,\dd\varphi
=2^{\alpha+1}K_+\,N^{-p}\int_{0}^{\infty}e^{-y^2}y^\alpha\,\dd y+o(N^{-p}).
\]
同理负侧给出 \(K_-\) 的同型项。最后利用
\(
\int_0^\infty e^{-y^2}y^\alpha\,\dd y=\frac12\Gamma\!\bigl(\frac{\alpha+1}{2}\bigr)=\frac12\Gamma(p)
\)
即可得到系数 \(c=2^\alpha\Gamma(p)(K_++K_-)\)。证毕。
\end{proof}

\begin{corollary}[dyadic 三级差分比值读出幂律指数]\label{cor:xi-endpoint-heat-probe-dyadic-exponent-readout}
在定理 \ref{thm:xi-endpoint-heat-probe-powerlaw-leading} 的假设下，令
\[
R_N:=\frac{a_N-a_{2N}}{a_{2N}-a_{4N}}.
\]
则
\begin{equation}\label{eq:xi-endpoint-heat-probe-dyadic-ratio-limit}
\boxed{\;\lim_{N\to\infty}R_N=2^{p}=2^{(\alpha+1)/2},\;}
\end{equation}
从而
\begin{equation}\label{eq:xi-endpoint-heat-probe-alpha-from-dyadic}
\boxed{\;\alpha=2\log_2\!\Bigl(\lim_{N\to\infty}R_N\Bigr)-1.\;}
\end{equation}
\end{corollary}
\begin{proof}
由 \eqref{eq:xi-endpoint-heat-probe-powerlaw-asymptotic}，
\(
a_N-a_{2N}=(a_N-m_{-1})-(a_{2N}-m_{-1})
=c(1-2^{-p})N^{-p}+o(N^{-p})
\)，
同理
\(
a_{2N}-a_{4N}=c(1-2^{-p})2^{-p}N^{-p}+o(N^{-p})
\)。
相除即得 \eqref{eq:xi-endpoint-heat-probe-dyadic-ratio-limit}，再反解得 \eqref{eq:xi-endpoint-heat-probe-alpha-from-dyadic}。证毕。
\end{proof}

\begin{corollary}[端点幂律系数的闭式核验提取]\label{cor:xi-endpoint-heat-probe-powerlaw-coefficient-extraction}
在定理 \ref{thm:xi-endpoint-heat-probe-powerlaw-leading} 的假设下，令
\[
D_N:=N^{p}(a_N-a_{2N}).
\]
则
\[
\lim_{N\to\infty}D_N=c(1-2^{-p}),
\]
从而
\begin{equation}\label{eq:xi-endpoint-heat-probe-Ksum-extraction}
\boxed{\;
K_++K_-=\frac{1}{2^\alpha\Gamma(p)}\cdot
\frac{\lim_{N\to\infty}D_N}{1-2^{-p}}\;}
\end{equation}
（其中 \(p=\tfrac{\alpha+1}{2}\)）。
\end{corollary}
\begin{proof}
由上一个推论的渐近式乘以 \(N^p\) 取极限即可，最后代入 \(c=2^\alpha\Gamma(p)(K_++K_-)\) 得 \eqref{eq:xi-endpoint-heat-probe-Ksum-extraction}。证毕。
\end{proof}

\paragraph{端点原子质量的外推抽取}
\begin{proposition}[Richardson 外推提升端点原子收敛阶]\label{prop:xi-endpoint-heat-probe-richardson-extrapolation}
设存在 \(p>0\) 与常数 \(c\neq 0\) 使
\(
a_N=m_{-1}+cN^{-p}+o(N^{-p})
\)。
定义
\[
\widehat m_N^{(p)}:=\frac{2^{p}a_{2N}-a_N}{2^{p}-1}.
\]
则 \(\widehat m_N^{(p)}=m_{-1}+o(N^{-p})\)。
若进一步有
\(
a_N=m_{-1}+cN^{-p}+c_1N^{-p-1}+o(N^{-p-1})
\)，则 \(\widehat m_N^{(p)}=m_{-1}+O(N^{-p-1})\)。
\end{proposition}
\begin{proof}
由假设
\(
a_{2N}=m_{-1}+c2^{-p}N^{-p}+o(N^{-p})
\)。
代入定义得
\(
2^pa_{2N}-a_N=(2^p-1)m_{-1}+o(N^{-p})
\)，除以 \(2^p-1\) 即得第一条。
含 \(N^{-p-1}\) 次项时同理直接抵消主项并得到阶提升。证毕。
\end{proof}

\begin{corollary}[自适应外推：无需预先给定 \(p\)]\label{cor:xi-endpoint-heat-probe-adaptive-extrapolation}
在定理 \ref{thm:xi-endpoint-heat-probe-powerlaw-leading} 的假设下，令
\[
\widehat p_N:=\log_2\!\Bigl(\frac{a_N-a_{2N}}{a_{2N}-a_{4N}}\Bigr),
\qquad
\widetilde m_N:=\frac{2^{\widehat p_N}a_{2N}-a_N}{2^{\widehat p_N}-1}.
\]
则 \(\widehat p_N\to p\) 且 \(\widetilde m_N\to m_{-1}\)。
\end{corollary}
\begin{proof}
由推论 \ref{cor:xi-endpoint-heat-probe-dyadic-exponent-readout} 得 \(\widehat p_N\to p\)。
又由命题 \ref{prop:xi-endpoint-heat-probe-richardson-extrapolation} 的公式对 \(p\) 的连续依赖性以及 \(a_N\to m_{-1}\) 得 \(\widetilde m_N\to m_{-1}\)。证毕。
\end{proof}

\paragraph{闭式卷积、生成函数与差分主项}
\begin{proposition}[中心二项式卷积闭式：线性时间计算]\label{prop:xi-endpoint-heat-probe-binomial-convolution}
对 Toeplitz 系数 \(\{c_h\}\)（\(|h|\le N\)）有闭式表示
\begin{equation}\label{eq:xi-endpoint-heat-probe-binomial-convolution}
\boxed{\;
a_N=\frac{1}{4^N}\sum_{h=-N}^{N}(-1)^h\binom{2N}{N+h}\,c_h.\;}
\end{equation}
因此 \(a_N\) 仅依赖 \(\{c_h\}_{|h|\le N}\)，并可由一次卷积在 \(O(N)\) 时间内计算。
\end{proposition}
\begin{proof}
由 \eqref{eq:xi-endpoint-heat-probe-bN-aN} 展开
\[
(b^{(N)})^\ast\mathbf{T}_N b^{(N)}
=\sum_{h=-N}^{N}c_h\sum_{k} \overline{b^{(N)}_{k+h}}\,b^{(N)}_{k},
\]
内和取使指标合法的 \(k\)。
由 \(b^{(N)}_{k}=(-1)^k\binom{N}{k}\) 得
\(
\overline{b^{(N)}_{k+h}}\,b^{(N)}_{k}=(-1)^h\binom{N}{k+h}\binom{N}{k}
\)。
对 \(k\) 求和并用 Vandermonde 卷积恒等式
\(
\sum_k\binom{N}{k+h}\binom{N}{k}=\binom{2N}{N+h}
\)
即得 \eqref{eq:xi-endpoint-heat-probe-binomial-convolution}。证毕。
\end{proof}

\begin{proposition}[端点母函数的 Stieltjes 表示与 Abel 残差读出]\label{prop:xi-endpoint-heat-probe-stieltjes-abel-residue}
令 \(\nu:=r_\#(\mu)\) 为 \eqref{eq:xi-endpoint-rxi-def} 的推前测度，则 \(\nu\) 支撑于 \([0,1]\) 且 \(\nu(\{1\})=m_{-1}\)。
并且对 \(N\ge 0\) 有
\[
a_N=\int_{[0,1]}x^N\,\dd\nu(x).
\]
定义生成函数
\(
\mathcal{A}(z):=\sum_{N\ge 0}a_N z^N
\)，则当 \(|z|<1\) 时有严格表示
\begin{equation}\label{eq:xi-endpoint-heat-probe-stieltjes-genfun}
\boxed{\;
\mathcal{A}(z)=\int_{[0,1]}\frac{1}{1-zx}\,\dd\nu(x).\;}
\end{equation}
并且端点原子质量满足 Abel 残差公式
\begin{equation}\label{eq:xi-endpoint-heat-probe-abel-residue}
\boxed{\;m_{-1}=\lim_{z\uparrow 1}(1-z)\mathcal{A}(z).\;}
\end{equation}
\end{proposition}
\begin{proof}
由 \eqref{eq:xi-endpoint-heat-probe-identity} 立得 \(a_N=\int x^N\,\dd\nu(x)\)。
对 \(|z|<1\) 有 \(\sum_{N\ge 0}(zx)^N=\frac{1}{1-zx}\) 且被 \(1/(1-|z|)\) 一致控制，可交换求和与积分得到 \eqref{eq:xi-endpoint-heat-probe-stieltjes-genfun}。
若 \(\nu\) 在 \(x=1\) 处有原子 \(m_{-1}\)，则 \(\mathcal{A}(z)=\frac{m_{-1}}{1-z}+(\text{在 }z=1\text{ 解析的项})\)，故 \((1-z)\mathcal{A}(z)\to m_{-1}\)。证毕。
\end{proof}

\begin{theorem}[端点探针的有限 Hankel 秩与有限共轭对刚性]\label{thm:xi-endpoint-probe-finite-hankel-rank-finite-conjugate-pairs}
在标量口径 \(\mathcal{A}=\CC\) 下，令 \(\mu\) 为 \(\TT\) 上有限正测度并定义端点探针
\[
a_N:=\frac{1}{4^N}\int_{\TT}|1-\xi|^{2N}\,\dd\mu(\xi),
\qquad
m_{-1}:=\mu(\{-1\}).
\]
令 \(\nu=r_\#(\mu)\)（\(r(\xi)=|1-\xi|^2/4\)），并记 \(s_N:=a_N-m_{-1}=\int_{[0,1)}x^N\,\dd\nu(x)\)。
则以下条件等价：
\begin{enumerate}
  \item 序列 \((s_N)_{N\ge 0}\) 的 Hankel 秩有限：存在 \(R<\infty\) 使得对一切 \(M\ge 0\)，Hankel 矩阵 \((s_{i+j})_{0\le i,j\le M}\) 的秩不超过 \(R\)；
  \item 测度 \(\nu\big|_{[0,1)}\) 为至多 \(R\) 个点的原子测度。
\end{enumerate}
在此情形下，存在互异 \(x_1,\dots,x_R\in(0,1)\) 与权重 \(w_j>0\) 使
\[
s_N=\sum_{j=1}^{R}w_j x_j^N,
\]
并且令
\(
Q(t):=\prod_{j=1}^{R}(t-x_j)=\sum_{k=0}^{R}q_k t^k
\)，则 \((s_N)\) 满足严格线性递推
\[
\sum_{k=0}^{R}q_k\,s_{N+k}=0\qquad(\forall N\ge 0).
\]
进一步，\(\mu\) 在 \(\TT\setminus\{-1\}\) 上的支撑包含于有限多个共轭点对 \(\{e^{\mathrm{i}\theta_j},e^{-\mathrm{i}\theta_j}\}\)（\(1\le j\le R\)），其中 \(x_j=\frac12(1-\cos\theta_j)\)。
\end{theorem}
\begin{proof}
由命题 \ref{prop:xi-endpoint-heat-probe-stieltjes-abel-residue}，\(s_N\) 为 \([0,1)\) 上正测度 \(\nu\big|_{[0,1)}\) 的 Hausdorff 矩序列。
Hausdorff 矩问题的结构理论给出：\((s_N)\) 的 Hankel 秩有限当且仅当该测度为有限原子；当
\(
\nu\big|_{[0,1)}=\sum_{j=1}^{R}w_j\delta_{x_j}
\)
时，必有 \(s_N=\sum_j w_jx_j^N\)。
此时递推由最小多项式 \(Q\) 直接给出：对任意 \(N\ge 0\)，
\[
\sum_{k=0}^{R}q_k s_{N+k}
=\sum_{j=1}^{R}w_j x_j^N\sum_{k=0}^{R}q_k x_j^k
=\sum_{j=1}^{R}w_j x_j^N Q(x_j)=0.
\]
最后，对任意 \(x\in(0,1)\)，方程 \(r(\xi)=x\) 在 \(\TT\) 上等价于 \(\Re\xi=1-2x\)，其解为一对共轭点 \(e^{\pm\mathrm{i}\theta}\)（\(\cos\theta=1-2x\)）；因此 \(\nu\big|_{[0,1)}\) 的每个原子对应 \(\TT\setminus\{-1\}\) 上的一个有限共轭对支撑。证毕。
\end{proof}

\begin{proposition}[幂律情形的纯差分主项]\label{prop:xi-endpoint-heat-probe-powerlaw-differences}
在定理 \ref{thm:xi-endpoint-heat-probe-powerlaw-leading} 的幂律假设下（\(p=\tfrac{\alpha+1}{2}\)），有
\[
a_N-a_{N+1}\sim p\,c\,N^{-p-1},
\qquad
a_{N-1}-2a_N+a_{N+1}\sim p(p+1)\,c\,N^{-p-2},
\]
其中 \(c\) 由 \eqref{eq:xi-endpoint-heat-probe-powerlaw-asymptotic} 给出。
\end{proposition}
\begin{proof}
由 \(a_N-m_{-1}=cN^{-p}+o(N^{-p})\)，令 \(u_N:=cN^{-p}\)。
则
\[
a_N-a_{N+1}=(u_N-u_{N+1})+o(N^{-p})-o(N^{-p})
=u_N\Bigl(1-(1+1/N)^{-p}\Bigr)+o(N^{-p-1}),
\]
而 \((1+1/N)^{-p}=1-\frac{p}{N}+O(N^{-2})\) 给出 \(u_N-u_{N+1}=pc\,N^{-p-1}+O(N^{-p-2})\)。
二阶中心差分同理由 \(n\mapsto n^{-p}\) 的二阶泰勒展开得到主项 \(p(p+1)c\,N^{-p-2}\)。
吸收误差即得。证毕。
\end{proof}

\endinput

